% 02_abstract.tex — Resumen en párrafo único (sin subtítulos IMRAD), conforme
% al formato de Scientia Agropecuaria: máx. 250 palabras (tolerancia 260).
% El contenido sigue el orden IMRAD: contexto/objetivo → métodos → resultados →
% conclusión. Las cifras provienen del embudo verificado (541→353→239→230→2);
% los valores aún no disponibles se definen en 00_datos_pendientes.tex.
\clearpage

\begin{abstract}

\noindent
El pH del jugo alcalinizado es la variable crítica del encalado azucarero,
y la ganancia variable de su curva de titulación degrada el control PID
convencional. El aprendizaje automático (AA) y el control avanzado se han
desarrollado en otros procesos de neutralización, pero su evidencia en el
proceso azucarero no estaba cartografiada. Se realizó una revisión de
alcance (PRISMA-ScR) en cinco bases (Scopus, WoS, IEEE~Xplore, SciELO,
LILACS; 2015--2026) con criterios \textit{a~priori}:
artículos de revista con cuartil SJR que aplicaran AA, control avanzado o
sensores virtuales al pH de la clarificación. El cribado fue doble,
ciego e independiente de los 230 registros por dos revisores, con
$\kappa$ inter-humano como métrica primaria ($\kappa=$\KappaInter). De 541 registros iniciales, dos estudios cumplieron
todos los criterios: un NMPC con identificación Takagi--Sugeno y la
sintonía metaheurística de un PID identificado con datos de planta;
ambos validaron el control solo en simulación. Un análisis de
sensibilidad halló otros 11 trabajos \textit{in-domain} excluidos por
criterios formales (3 artículos en revistas sin cuartil y 8 actas),
también basados en simulación. Ninguno de los 13 aplicó aprendizaje por
refuerzo ni propone un sensor virtual de pH en línea, ninguno declara
control en lazo cerrado en planta industrial y 10 provienen de dos
grupos de autores. La brecha es de madurez y diversidad metodológica,
no de ausencia de investigación. Se propone una agenda por horizontes:
control neuro-difuso y sensores virtuales, NMPC de caja gris y
aprendizaje por refuerzo con validación en planta.

\medskip
\keywords{control de pH; industria azucarera; aprendizaje automático;
control predictivo; revisión de alcance; análisis bibliométrico;
brecha de investigación}

\end{abstract}
\clearpage

% ── Abstract en inglés (obligatorio según la guía de Scientia Agropecuaria,
%    incluso cuando el cuerpo del manuscrito está en español) ────────────────
\begin{otherlanguage}{english}
\begin{abstract}

\noindent
The pH of alkalinized juice is the critical variable of sugar-industry
liming, and the variable gain of its titration curve degrades
conventional PID control. Machine learning (ML) and advanced control have
been developed for other neutralization processes, but their evidence in
the sugar process had not been mapped. A scoping review (PRISMA-ScR) was
conducted in Scopus, Web of Science, IEEE~Xplore, SciELO and LILACS
(2015--2026) with \textit{a~priori} criteria: journal articles with an
SJR quartile (Q1--Q4) applying ML, advanced control or virtual sensors to
clarification or liming pH. Screening was double, blinded and
independent, by two reviewers on all 230 records, with inter-rater
$\kappa$ as the primary metric ($\kappa=$\KappaInter). Of 541 initial records, two studies met all criteria: an
NMPC with Takagi--Sugeno identification, and metaheuristic PID tuning
with identification from plant data; control was validated only in
simulation. A sensitivity analysis found 11 further in-domain works
excluded by formal criteria (3 articles in journals without a quartile
and 8 conference papers or chapters), also simulation-based. None of the
13 applied reinforcement learning or proposes a virtual pH sensor for
online use, and none reports closed-loop control in an industrial plant;
10 come from two author groups. The gap concerns maturity and methodological diversity, not an
absence of research. A horizon-based agenda is proposed: neuro-fuzzy
control and virtual sensors, grey-box NMPC, and reinforcement learning
with plant validation.

\medskip
\keywords{pH control; sugar industry; machine learning; predictive
control; scoping review; bibliometric analysis; research gap}

\end{abstract}
\end{otherlanguage}
\clearpage
